% profiles/university-of-missouri-columbia.tex — University of Missouri-Columbia
% ============================================================================
% source: https://gradschool.missouri.edu/academics/thesis-dissertation/formatting-guidelines/
%         (Graduate School, "Formatting Your Thesis or Dissertation" — the
%         section-by-section guide, with links to a specimen of every page)
% source: https://gradschool.missouri.edu/policy/body-format-for-thesis-or-dissertation/
%         (Graduate School policy, "Body Format for Thesis or Dissertation" —
%         the same figures stated as policy)
% source: https://gradschool.missouri.edu/policy/electronic-dissertation-and-thesis-basics/
% source: https://gradschool.missouri.edu/academics/thesis-dissertation/technology/
%         (both read for accessibility, which neither mentions)
% source: https://gradschool.missouri.edu/wp-content/uploads/2025/05/Example-Dissertation-Title-Page.pdf
% source: https://gradschool.missouri.edu/wp-content/uploads/2025/05/Example-thesis-title-page.pdf
% source: https://gradschool.missouri.edu/wp-content/uploads/2025/05/approval-page.pdf
%         (the three specimens the guide links; uploaded 2025-05)
% accessed: 2026-09-12
% checked:  texlib-thesis-profile-round
%
% VERIFIED and set: geometry, spacing, the title page and the approval page.
% Accessibility is recorded as an absence — see below.
%
% NOT set, for want of a hook: the font rules. "The font used should be a
% commonly available font, such as Arial, Veranda [sic], Times New Roman, etc.";
% "Use 10-point or 12-point fonts for regular text. Headings may be 14-point or
% 16-point. Tables, figures, etc., may be smaller than 10-point." The list is
% open ("etc."), so Latin Modern is not excluded, but it is not named either.
%
% Mizzou's two sources say the same thing in the same words, which is unusually
% convenient: the policy page and the formatting guide agree on every figure
% here, so there is no prose-versus-specimen conflict to resolve.

\thesisinstitution{University of Missouri-Columbia}

% "Margins: Left Margin – 1.5 inches; Top, Bottom, and Right Margins – 1 inch",
% stated identically on the policy page and the formatting guide. No reason is
% given for the wider left margin and no exception is offered; it applies to the
% filed research.pdf. Being a LEFT margin rather than an inner one, a twoside
% build would put it on the wrong side of even pages — keep Mizzou documents
% oneside.
\thesissetgeometry{left=1.5in, right=1in, top=1in, bottom=1in}

% "Line Spacing: All standard text should be double-spaced. References,
% bibliographic works, and endnotes should be consistent with the formatting
% standards of your selected style manual." The front matter has its own rules
% (the List of Illustrations is "Single-space lines within entries, and
% double-space between entries"), which are not profile fields.
% \thesissetchapteropening is NOT called: no source mentions one.
\thesissetspacing{double}

% --- Accessibility -----------------------------------------------------------
% NOT PUBLISHED. The formatting guide, the two policy pages and the "Technology
% & Your Submission" page were all read on the access date and none of them
% mentions accessibility, tagging, alt text, WCAG, PDF/A or the ADA. The
% Technology page is where such a rule would sit — it is the page about file
% formats and PDF creation — and it goes from "All thesis and dissertation
% materials must be submitted in pdf format" straight to multimedia formats and
% copyright.
%
% So no \thesisrequiretagging, no \thesisrequiredocumentlanguage, no
% \thesisrequirepdfstandard. This class emits a tagged PDF regardless; declaring
% that Mizzou asked for one would be inventing a requirement.
\thesisaccessibilityrequirement{Not published}
	{The Graduate School's filing requirements state no accessibility rule. The
	 only related instruction is the file format itself: "All thesis and
	 dissertation materials must be submitted in pdf format", in a single file
	 named exactly "research.pdf", all lowercase.}

% --- Title page --------------------------------------------------------------
% "The body of the text should be centered on the page, both vertically and
% horizontally", the page is counted as Roman numeral i but carries no number,
% and — a trap worth naming — "Use the month and year of your graduation, not
% the month and year of your defense." So \graddate is the GRADUATION date here,
% and the specimen sets it in capitals ("MAY 2001").
%
% The block below is the specimen, which is the same page for a thesis and a
% dissertation with the noun swapped: the title in capitals, a rule, the
% "presented to" block, a rule, the partial-fulfilment block, the degree, a
% rule, then "by", the author in capitals, and the supervisor line.
%
% The specimen draws its three rules at three slightly different widths. That
% reads as an artifact of the Word original rather than a rule — nothing in the
% prose mentions the rules at all — so one consistent width is used here.
%
% The supervisor line is "Dr. Thomas Sink, Dissertation Supervisor", i.e. the
% name exactly as \gradadvisor gives it, followed by the role. Mizzou's specimen
% writes the courtesy title as "Dr."; the class makes no assumption, so put
% whatever form is wanted into \gradadvisor.
\renewcommand{\thesistitlepage}{%
	\loadgeometry{nopagenumber}%
	\begin{titlepage}%
		\thispagestyle{empty}\singlespacing\centering
		\vspace*{\fill}
		\thesis@tagtitle{Title}{\MakeUppercase{\@title}}\par
		\vspace{2\baselineskip}
		\rule{0.6\textwidth}{0.4pt}\par
		\vspace{2\baselineskip}
		A \thesis@DocNoun\\
		presented to\\
		the Faculty of the Graduate School\\
		at the University of Missouri-Columbia\par
		\vspace{2\baselineskip}
		\rule{0.6\textwidth}{0.4pt}\par
		\vspace{2\baselineskip}
		In Partial Fulfillment\\
		of the Requirements for the Degree\par
		\vspace{2\baselineskip}
		\thesis@degree\par
		\vspace{2\baselineskip}
		\rule{0.6\textwidth}{0.4pt}\par
		\vspace{2\baselineskip}
		by\par
		\MakeUppercase{\@author}\par
		\thesis@advisor, \thesis@DocNoun{} Supervisor\par
		\vspace{2\baselineskip}
		\MakeUppercase{\thesis@date}\par
		\vspace*{\fill}
	\end{titlepage}%
	\loadgeometry{normal}%
}

% --- Committee approval page --------------------------------------------------
% Mizzou REQUIRES one, and it is the third page of research.pdf — after the
% title page, and after the optional copyright page if there is one. Its rules:
%
%   "The approval page is not counted, numbered, or listed in the Table of
%    Contents." / "The title must be centered." / "Each person on your committee
%    must be listed." / "The approval page in your thesis/dissertation should
%    not have committee signatures."
%
% That last line is the important one: NO SIGNATURE RULES. Signatures are
% collected outside the document; the filed page carries typed names only. The
% class's default member line draws a rule, so it is replaced here.
%
% The wording below is the specimen's, with its bracketed placeholders filled
% from the class's metadata. The specimen's opening paragraph is flush left and
% the title is centred, exactly as the prose requires.
%
% ONE THING TO KNOW ABOUT THE NAMES: the specimen prefixes each with
% "Professor" — "Professor Jane Doe", "Professor Vijay Kumar". This file does
% NOT add that prefix, because it would be wrong for a committee member who is
% not a professor, and the class cannot tell. Put the form you want into
% \committeemember, e.g. \committeemember{Professor Jane Doe}{}.
%
% The specimen also writes the degree in lower case inside its placeholder
% ("a candidate for the degree of [doctor of philosophy, doctor of education or
% master of ________]"). Since that is a placeholder listing the options rather
% than a typographic instruction, \thesis@degree is printed as the filer set it.
\renewcommand{\thesisapprovalpage}{%
	\clearpage\thispagestyle{empty}%
	\loadgeometry{nopagenumber}\singlespacing
	\noindent The undersigned, appointed by the dean of the Graduate School, have
	examined the \thesis@docnoun{} entitled\par
	\vspace{\baselineskip}
	\begin{center}
		\MakeUppercase{\@title}
	\end{center}
	\vspace{\baselineskip}
	\noindent presented by \@author,\\
	a candidate for the degree of \thesis@degree,\\
	and hereby certify that, in their opinion, it is worthy of acceptance.\par
	\vspace{3\baselineskip}
	\begin{center}
		\begingroup
		% No signature rules: "The approval page in your thesis/dissertation
		% should not have committee signatures."
		\renewcommand{\thesis@memberline}[2]{##1\par\vspace{\baselineskip}}%
		\thesis@committee
		\endgroup
	\end{center}%
	\loadgeometry{normal}\doublespacing\clearpage
}
